\chapter{Четырёхкомпонентное происхождение ковекторов источника KMS}
\label{ch:v9-endpoint-kms-source-covector-origin}

\section{Две масштабные связи}

Пусть $s,a,t$ обозначают два singlet-типа и charged family-triplet.
Multiplicity trace задаёт веса
\begin{equation}
 w=(1,1,3)^T,
 \qquad W=\operatorname{diag}(1,1,3).
 \label{eq:v9-kms-covector-origin-weights}
\end{equation}
Из четырёх slots непрерывная пара $(E_*,\chi)$ задаёт ровно два общих
размерных масштаба:
\begin{equation}
 \Delta_*=E_*,
 \qquad
 \Gamma_*=\chi^2\frac{E_*}{\hbar}.
 \label{eq:v9-kms-covector-origin-common-scales}
\end{equation}
В логарифмических координатах соответствующая карта имеет якобиан
\begin{equation}
 J_{\rm scale}=
 \frac{\partial(\log\Delta_*,\log\Gamma_*)}
 {\partial(\log E_*,\log\chi)}
 =\begin{pmatrix}1&0\\1&2\end{pmatrix},
 \qquad \rank J_{\rm scale}=2,
 \qquad \det J_{\rm scale}=2.
 \label{eq:v9-kms-covector-origin-scale-jacobian}
\end{equation}
Следовательно, общий gap-scale и общий conductance-scale действительно
связаны с уже имеющимися slots без нового размерного параметра.

\section{Разложение на scale и shape}

Запишем три gaps и три conductances как
\begin{equation}
 \boldsymbol\Delta=E_*r_\theta,
 \qquad
 \boldsymbol\kappa=\Gamma_*r_\kappa,
 \qquad r_\theta,r_\kappa\in\mathbb R_{>0}^3,
 \label{eq:v9-kms-covector-origin-factorization}
\end{equation}
и удалим повторный общий масштаб двумя weighted-нормировками
\begin{equation}
 w^Tr_\theta=5,
 \qquad w^Tr_\kappa=5.
 \label{eq:v9-kms-covector-origin-normalizations}
\end{equation}
Каждая положительная shape-поверхность двумерна. Удобный tangent-базис
общего ядра имеет вид
\begin{equation}
 B=\begin{pmatrix}
 1&0\\0&1\\-1/3&-1/3
 \end{pmatrix},
 \qquad w^TB=0,
 \qquad \rank B=2.
 \label{eq:v9-kms-covector-origin-shape-basis}
\end{equation}
Тем самым после фиксации двух scales остаются четыре независимые
относительные координаты, например
\begin{equation}
 \rho_{\theta s}=\frac{\Delta_s}{\Delta_t},\quad
 \rho_{\theta a}=\frac{\Delta_a}{\Delta_t},\quad
 \rho_{\kappa s}=\frac{\kappa_s}{\kappa_t},\quad
 \rho_{\kappa a}=\frac{\kappa_a}{\kappa_t}.
 \label{eq:v9-kms-covector-origin-relative-ratios}
\end{equation}
После выбора единиц соответствующие covectors предыдущего гейта получаются
метрическим двойствованием
\begin{equation}
 j_\theta=W\boldsymbol\Delta,
 \qquad
 j_\kappa=W\boldsymbol\kappa.
 \label{eq:v9-kms-covector-origin-metric-duals}
\end{equation}

\section{Полный rank-аудит}

Параметризуем нормированные shapes формулами
\begin{equation}
 r_\theta(a,b)=\left(a,b,\frac{5-a-b}{3}\right)^T,
 \qquad
 r_\kappa(c,d)=\left(c,d,\frac{5-c-d}{3}\right)^T.
 \label{eq:v9-kms-covector-origin-shape-chart}
\end{equation}
Для карты
\begin{equation}
 \mathcal F(E_*,\chi,a,b,c,d)
 =\left(E_*r_\theta(a,b),
 \chi^2\frac{E_*}{\hbar}r_\kappa(c,d)\right)
 \label{eq:v9-kms-covector-origin-full-map}
\end{equation}
при $\hbar=E_*=\chi=a=b=c=d=1$ машинный аудит даёт
\begin{equation}
 \rank D\mathcal F=6,
 \qquad \det D\mathcal F=\frac{50}{9}.
 \label{eq:v9-kms-covector-origin-full-rank}
\end{equation}
Два scale-столбца имеют ранг $2$, четыре shape-столбца --- ранг $4$:
\begin{equation}
 \rank D\mathcal F_{\rm scale}=2,
 \qquad
 \rank D\mathcal F_{\rm shape}=4,
 \qquad
 6-2=4.
 \label{eq:v9-kms-covector-origin-rank-defect}
\end{equation}

\section{Endpoint types и transport}

Endpoint-extension фиксирует само разложение типов $1\oplus1\oplus3$ и
потому различает три координатные линии. Но rank-three type-system не
задаёт точку на двух shape-поверхностях. Transport может лишь выбрать
ориентацию уже построенных каналов. Для двух допустимых представителей
\begin{equation}
 P_+=I_3,
 \qquad
 P_-=\begin{pmatrix}0&1&0\\1&0&0\\0&0&1\end{pmatrix}
 \label{eq:v9-kms-covector-origin-transport-representatives}
\end{equation}
выполнено
\begin{equation}
 P_\pm^TWP_\pm=W,
 \qquad
 \rank(P_\pm B)=2.
 \label{eq:v9-kms-covector-origin-transport-isometry}
\end{equation}
Поэтому дискретная ориентация не уменьшает непрерывную shape-размерность.

\section{Точные свидетели неединственности}

При одинаковых $E_*=\Gamma_*=1$ рассмотрим
\begin{equation}
 r_\theta^{(A)}=(1,1,1)^T,
 \qquad
 r_\theta^{(B)}=(2,1,2/3)^T,
 \label{eq:v9-kms-covector-origin-gap-witnesses}
\end{equation}
и
\begin{equation}
 r_\kappa^{(A)}=(1,1,1)^T,
 \qquad
 r_\kappa^{(B)}=(1,2,2/3)^T.
 \label{eq:v9-kms-covector-origin-conductance-witnesses}
\end{equation}
Все четыре вектора положительны и имеют weighted-нормировку $5$, но дают
разные covectors:
\begin{equation}
 Wr_\theta^{(A)}\ne Wr_\theta^{(B)},
 \qquad
 Wr_\kappa^{(A)}\ne Wr_\kappa^{(B)}.
 \label{eq:v9-kms-covector-origin-distinct-witnesses}
\end{equation}
Одинаковые four-slot scales, endpoint types и transport orientation тем
самым совместимы с различными KMS packages.

\section{Origin-ledger}

\begin{equation}
 \begin{aligned}
 N_{\rm common\ scale\ links}&=2/2,\\
 N_{\rm channel\ type\ labels}&=3/3,\\
 N_{\rm relative\ ratio\ origin}&=0/4,\\
 N_{\rm source\ covector\ origin}&=0/2,\\
 N_{\rm physical\ four\ slot\ parent}&=0/1.
 \end{aligned}
 \label{eq:v9-kms-covector-origin-ledger}
\end{equation}

\begin{theorem}[Four-slot factorization ковектор источникаs]
Четыре slots фиксируют общий масштаб gaps $E_*$, общий масштаб
conductances $\chi^2E_*/\hbar$, разложение трёх channel types и допустимую
ориентацию transport. Однако полный six-component source-package имеет
четыре дополнительные относительные shape-координаты. Поэтому
$j_\theta,j_\kappa$ не выводятся из четырёх slots однозначно.
\end{theorem}

\begin{nogocriterion}[Тип и ориентация не являются численным selector]
Различимость трёх channel types и выбор transport orientation не создают
четырёх скалярных уравнений для
$(\rho_{\theta s},\rho_{\theta a},\rho_{\kappa s},\rho_{\kappa a})$.
Без отдельного bounded selector этих величин ковектор источникаs остаются
внешними данными общего родительского функционала.
\end{nogocriterion}

\begin{protocol}\begin{enumerate}
\item общий gap-scale: \textbf{$E_*$};
\item общий conductance-scale: \textbf{$\chi^2E_*/\hbar$};
\item scale-map: \textbf{rank $2$, determinant $2$};
\item channel types: \textbf{$3/3$};
\item transport orientation: \textbf{shape-rank сохраняется};
\item свободная relative dimension: \textbf{$4$};
\item relative-ratio origin: \textbf{$0/4$};
\item ковектор источника origin: \textbf{$0/2$};
\item следующий гейт: \textbf{minimal selector четырёх относительные формы}.
\end{enumerate}\end{protocol}

Машинный сертификат сохранён в
\path{s2t_v9_endpoint_creation_kms_source_covector_four_slot_parent_origin_gate_results.json}.